\textbf{Phần 1. Trắc nghiệm lựa chọn}
\setcounterpageref{ex}{0}

\begin{ex} %[3H3N17]
Đoạn thẳng nối tâm với một điểm bất kì trên đường tròn được gọi là
\choice
{\True Bán kính}
{Đường kính}
{Dây cung}
{Chu vi}
\loigiai{
Đoạn thẳng nối tâm với một điểm trên đường tròn là bán kính.\\
\textbf{Đáp án của câu hỏi là A}
}
\end{ex}

\begin{ex} %[3H3N17]
Đoạn thẳng đi qua tâm và nối hai điểm trên đường tròn được gọi là
\choice
{Bán kính}
{Tâm}
{\True Đường kính}
{Cung tròn}
\loigiai{
Đoạn thẳng đi qua tâm nối hai điểm trên đường tròn là đường kính.\\
\textbf{Đáp án của câu hỏi là C}
}
\end{ex}

\begin{ex} %[3H3H17]
Độ dài đường kính gấp mấy lần độ dài bán kính của cùng một hình tròn?
\choice
{$3$ lần}
{\True $2$ lần}
{$4$ lần}
{$1$ lần}
\loigiai{
Đường kính gấp 2 lần bán kính ($d = 2 \times r$).\\
\textbf{Đáp án của câu hỏi là B}
}
\end{ex}

\begin{ex} %[3H3H17]
Một hình tròn có bán kính $4\text{ cm}$. Đường kính của hình tròn đó là
\choice
{$6\text{ cm}$}
{$10\text{ cm}$}
{\True $8\text{ cm}$}
{$16\text{ cm}$}
\loigiai{
$4 \times 2 = 8\text{ cm}$.\\
\textbf{Đáp án của câu hỏi là C}
}
\end{ex}

\begin{ex} %[3H3H17]
Một hình tròn có đường kính $12\text{ cm}$. Bán kính của hình tròn đó là
\choice
{$4\text{ cm}$}
{\True $6\text{ cm}$}
{$8\text{ cm}$}
{$24\text{ cm}$}
\loigiai{
$12 : 2 = 6\text{ cm}$.\\
\textbf{Đáp án của câu hỏi là B}
}
\end{ex}

\begin{ex} %[3H3H17]
Tâm $O$ của hình tròn là
\choice
{Điểm nằm trên đường tròn}
{Điểm nằm ngoài hình tròn}
{Đoạn thẳng}
{\True Trung điểm của mọi đường kính}
\loigiai{
Tâm $O$ là trung điểm của mọi đường kính.\\
\textbf{Đáp án của câu hỏi là D}
}
\end{ex}

\begin{ex} %[3H3V17]
Một hình tròn có bán kính $7\text{ cm}$. Độ dài đường kính của hình tròn là
\choice
{$12\text{ cm}$}
{$16\text{ cm}$}
{\True $14\text{ cm}$}
{$21\text{ cm}$}
\loigiai{
$7 \times 2 = 14\text{ cm}$.\\
\textbf{Đáp án của câu hỏi là C}
}
\end{ex}

\begin{ex} %[3H3V17]
Một hình tròn có đường kính $18\text{ dm}$. Bán kính của hình tròn đó là
\choice
{$8\text{ dm}$}
{$10\text{ dm}$}
{$36\text{ dm}$}
{\True $9\text{ dm}$}
\loigiai{
$18 : 2 = 9\text{ dm}$.\\
\textbf{Đáp án của câu hỏi là D}
}
\end{ex}

\begin{ex} %[3H3V17]
Dùng com-pa mở rộng 2 đầu kim cách nhau $5\text{ cm}$ để vẽ một hình tròn. Hình tròn vẽ được có đường kính là
\choice
{\True $10\text{ cm}$}
{$5\text{ cm}$}
{$15\text{ cm}$}
{$20\text{ cm}$}
\loigiai{
Độ mở com-pa là bán kính: $r = 5\text{ cm} \Rightarrow$ Đường kính: $5 \times 2 = 10\text{ cm}$.\\
\textbf{Đáp án của câu hỏi là A}
}
\end{ex}

\begin{ex} %[3H3V17]
Một miệng giếng hình tròn có bán kính $60\text{ cm}$. Đường kính miệng giếng đó là
\choice
{$100\text{ cm}$}
{$140\text{ cm}$}
{$30\text{ cm}$}
{\True $120\text{ cm}$}
\loigiai{
$60 \times 2 = 120\text{ cm}$.\\
\textbf{Đáp án của câu hỏi là D}
}
\end{ex}

\begin{ex} %[3H3V17]
Cho hình tròn tâm $O$ có đường kính $AB = 20\text{ cm}$. Độ dài đoạn thẳng $OA$ là
\choice
{$5\text{ cm}$}
{\True $10\text{ cm}$}
{$15\text{ cm}$}
{$20\text{ cm}$}
\loigiai{
$OA$ là bán kính: $20 : 2 = 10\text{ cm}$.\\
\textbf{Đáp án của câu hỏi là B}
}
\end{ex}

\begin{ex} %[3H3C17]
Cho hai hình tròn cùng tâm $O$: hình tròn nhỏ có bán kính $3\text{ cm}$, hình tròn lớn có bán kính $5\text{ cm}$. Khoảng cách từ một điểm trên đường tròn nhỏ đến một điểm trên đường tròn lớn cùng nằm trên một tia gốc $O$ là
\choice
{\True $2\text{ cm}$}
{$1\text{ cm}$}
{$3\text{ cm}$}
{$8\text{ cm}$}
\loigiai{
Khoảng cách: $5 - 3 = 2\text{ cm}$.\\
\textbf{Đáp án của câu hỏi là A}
}
\end{ex}

\textbf{Phần 2. Trắc nghiệm đúng sai}
\setcounterpageref{ex}{0}

\begin{ex} %[3H3H17]
Xét tính đúng sai về các yếu tố của hình tròn:
\choiceTF
{\True Trong một hình tròn, tất cả các bán kính đều có độ dài bằng nhau.}
{\True Trong một hình tròn, tất cả các đường kính đều có độ dài bằng nhau.}
{Độ dài bán kính dài gấp đôi độ dài đường kính.}
{\True Độ dài đường kính dài gấp đôi độ dài bán kính.}
\loigiai{
\begin{itemize}
    \item Các bán kính bằng nhau: Đúng.
    \item Các đường kính bằng nhau: Đúng.
    \item Đường kính gấp đôi bán kính (không phải bán kính gấp đôi đường kính). Ý 3 Sai, Ý 4 Đúng.
\end{itemize}
\begin{itemize}
    \item a) \textbf{Đúng}.
    \item b) \textbf{Đúng}.
    \item c) \textbf{Sai}.
    \item d) \textbf{Đúng}.
\end{itemize}
}
\end{ex}

\begin{ex} %[3H3V17]
Một hình tròn có đường kính $16\text{ cm}$. Xét tính đúng sai:
\choiceTF
{\True Bán kính của hình tròn là $8\text{ cm}$.}
{\True Tâm $O$ cách mọi điểm trên đường tròn một khoảng $8\text{ cm}$.}
{Nếu vẽ thêm đường kính thứ hai thì hai đường kính có độ dài khác nhau.}
{\True Mọi đường kính vẽ trong hình tròn này đều dài $16\text{ cm}$.}
\loigiai{
\begin{itemize}
    \item Bán kính: $16 : 2 = 8\text{ cm}$. Đúng.
    \item Khoảng cách từ tâm đến đường tròn là bán kính $= 8\text{ cm}$. Đúng.
    \item Các đường kính đều bằng nhau ($16\text{ cm}$). Ý 3 Sai, Ý 4 Đúng.
\end{itemize}
\begin{itemize}
    \item a) \textbf{Đúng}.
    \item b) \textbf{Đúng}.
    \item c) \textbf{Sai}.
    \item d) \textbf{Đúng}.
\end{itemize}
}
\end{ex}

\textbf{Phần 3. Trả lời ngắn}
\setcounterpageref{ex}{0}

\begin{ex} %[3H3H17]
Một hình tròn có bán kính $6\text{ cm}$. Tính đường kính của hình tròn theo xăng-ti-mét.
\shortans{12}
\loigiai{
\textbf{Đáp số:} $12$\\
$6 \times 2 = 12\text{ cm}$.
}
\end{ex}

\begin{ex} %[3H3H17]
Một hình tròn có đường kính $14\text{ cm}$. Tính bán kính của hình tròn theo xăng-ti-mét.
\shortans{7}
\loigiai{
\textbf{Đáp số:} $7$\\
$14 : 2 = 7\text{ cm}$.
}
\end{ex}

\begin{ex} %[3H3V17]
Để vẽ đường tròn có đường kính $16\text{ cm}$, ta cần mở com-pa một khoảng bằng bao nhiêu xăng-ti-mét?
\shortans{8}
\loigiai{
\textbf{Đáp số:} $8$\\
$16 : 2 = 8\text{ cm}$.
}
\end{ex}

\begin{ex} %[3H3V17]
Một bánh xe đạp có bán kính $30\text{ cm}$. Tính đường kính bánh xe đạp theo xăng-ti-mét.
\shortans{60}
\loigiai{
\textbf{Đáp số:} $60$\\
$30 \times 2 = 60\text{ cm}$.
}
\end{ex}

\begin{ex} %[3H3V17]
Một hình tròn có đường kính $20\text{ dm}$. Tính bán kính của hình tròn theo đề-xi-mét.
\shortans{10}
\loigiai{
\textbf{Đáp số:} $10$\\
$20 : 2 = 10\text{ dm}$.
}
\end{ex}

\begin{ex} %[3H3C17]
Cho đoạn thẳng $AB = 10\text{ cm}$. Vẽ hình tròn tâm $A$ bán kính $AB$. Đường kính của hình tròn đó dài bao nhiêu xăng-ti-mét?
\shortans{20}
\loigiai{
\textbf{Đáp số:} $20$\\
Bán kính $r = 10\text{ cm} \Rightarrow$ Đường kính: $10 \times 2 = 20\text{ cm}$.
}
\end{ex}
